% ---------------------------------------------------------------------
%  附录 E 测试矩阵
% ---------------------------------------------------------------------
\chapter{测试矩阵}
\label{app:tests}

\section{单元测试分布}
\label{sec:unit-distribution}

\begin{table}[H]
    \centering
    \caption[各模块的单元测试数]{各模块的单元测试数}
    \label{tab:unit-dist}
    \small
    \begin{tabular}{@{}lr@{}}
        \toprule
        模块 & 测试数 \\
        \midrule
        \code{complex}     & 5 \\
        \code{linalg}      & 26 \\
        \code{math}        & 26 \\
        \code{types}       & 27 \\
        \code{basics}      & 53 \\
        \code{weights}     & 8 \\
        \code{convolution} & 26 \\
        \code{dyson}       & 3 \\
        \code{vie2}        & 3 \\
        \code{equilibrium} & 9 \\
        \code{fourier}     & 49 \\
        \code{observables} & 2 \\
        \code{io}          & 15 \\
        \midrule
        合计 & 252 \\
        \bottomrule
    \end{tabular}
\end{table}

\section{集成测试分布}
\label{sec:integration-distribution}

\begin{table}[H]
    \centering
    \caption[集成测试目标]{集成测试目标}
    \label{tab:integration-dist}
    \small
    \begin{tabular}{@{}lrrl@{}}
        \toprule
        目标 & 文件数 & 测试数 & 主要内容 \\
        \midrule
        \code{bubble}      & 3 & 12 & 气泡图三分量 \\
        \code{convolution} & 4 & 20 & 卷积与增量变体 \\
        \code{dyson}       & 4 & 40 & Matsubara、启动、推进 \\
        \code{fourier}     & 5 & 20 & FFT、修正、谱函数 \\
        \code{io}          & 3 & 27 & 四种格式的往返与错误 \\
        \code{math}        & 3 & 5  & 数学工具与阶数门 \\
        \code{moving}      & 6 & 29 & 窗口、推进、截断解 \\
        \code{pseudo}      & 6 & 26 & 赝粒子全流程 \\
        \code{steady}      & 5 & 23 & 稳态 Dyson 与观测量 \\
        \code{types}       & 5 & 43 & 容器与操作 \\
        \code{weights}     & 2 & 4  & 权重与阶数 \\
        \code{vie}         & 1 & 22 & VIE2 全流程与奇异核 \\
        \code{config_parser} & 1 & 7 & 配置解析 \\
        \midrule
        合计 & 48 & 278 & \\
        \bottomrule
    \end{tabular}
\end{table}

\section{测试类型}
\label{sec:test-kinds}

\begin{table}[H]
    \centering
    \caption[测试类型与代表用例]{测试类型与代表用例}
    \label{tab:test-kinds}
    \small
    \begin{tabularx}{\textwidth}{@{}lX@{}}
        \toprule
        类型 & 代表用例 \\
        \midrule
        解析值 & 常数核的卷积等于区间长度；单极点的解析表达式 \\
        构造 & 容器的字段、块数、越界断言 \\
        一致性 & 定长块乘与通用块乘逐位相同；缓冲与对象目标相同 \\
        收敛 & 误差随 $h$ 的下降阶数 \\
        参考数据 & 与独立实现的逐点比较 \\
        错误输入 & 非法网格、重复索引、非有限值 \\
        往返 & 写入后读出，逐位相等 \\
        \bottomrule
    \end{tabularx}
\end{table}
